%% ===================== 7. KẾT LUẬN =====================
\section{Kết luận và Hướng phát triển}\label{sec:concl}
Báo cáo áp dụng loại hình Applied Experimentation để (i) tái lập một detector kiểu DeSFAM (VAE + Isolation Forest) cho anomaly detection trên system call với dữ liệu DongTing thực (AUC $0{,}832$ sau sensor-alphabet filter) và (ii) định lượng hệ quả an ninh của việc cập nhật ngưỡng \textit{vô điều kiện}: khi triển khai theo nghĩa đen công thức EMA của DeSFAM, noise từ một noisy neighbor làm inflate ngưỡng và giúp $74{,}2\%$ tấn công stealthy đạt evasion, với xu hướng thống kê rõ rệt theo cường độ nhiễu ($\rho=-0{,}74$). Ngược lại, dạng cập nhật \textit{có điều kiện} mà chính \textit{Algorithm 2} của DeSFAM đặc tả vô hiệu hoá hoàn toàn vector này ($0\%$). Qua đó, báo cáo vạch ra một sự thiếu nhất quán trong chính bài báo DeSFAM (công thức EMA vô điều kiện so với Algorithm 2 có điều kiện) và chứng minh rằng quy tắc cập nhật có điều kiện là \textit{thiết yếu về an ninh}, không chỉ để giảm báo động giả. Hướng phát triển tiếp theo gồm: (1) xác thực online trên cụm K8s + Tetragon; (2) khảo sát ngưỡng theo từng container; và (3) bổ sung argument analysis cho system call nhằm chống tấn công stealthy triệt để hơn.
